\documentclass[openany]{article}

%Typesetting and language
\usepackage[american]{babel}
\usepackage[T1]{fontenc}
\usepackage{charter}
\usepackage{enumitem}
\usepackage{hyperref}

%Symbols
\usepackage{amssymb, amsmath, amsthm, bm}
\usepackage{mathrsfs}
\usepackage{mathtools}
\usepackage{marvosym}
\usepackage{MnSymbol}

%Colors & graphics
\usepackage[dvipsnames]{xcolor}
\usepackage{pgfplots}
\usepackage[numbered,framed]{matlab-prettifier}
\usepackage{pgfplots}
\usepackage{listings}
\usepackage{tikz}
\usetikzlibrary{arrows.meta}
\usepackage[object=vectorian]{pgfornament}
\usepackage{wrapfig}
\usepackage{varwidth}
\usepackage[framemethod=TikZ]{mdframed}
\usepackage{caption}
\usepackage{float}
\usepackage{geometry}
\usepackage{ulem}
\usepackage[most]{tcolorbox}
\usepackage{array}

\setlength{\parindent}{0pt}

\makeatletter
\g@addto@macro\bfseries{\boldmath}
\makeatother


\renewcommand{\Re}{\mathfrak{Re}}
\renewcommand{\Im}{\mathfrak{Im}}

\geometry{left=2cm,right=2cm,bottom=2cm,top=2cm}

\usepackage{fancyhdr}
\pagestyle{fancy}
\fancyhf{}
\renewcommand{\sectionmark}[1]{\markright{\arabic{section} - #1}}
\cfoot{\thepage}
\lhead{Algo Learning}
\chead{Dynamic Programming}
\rhead{Adam Yang}
\renewcommand{\headrulewidth}{1pt}


\DeclareMathOperator{\sgn}{sgn}
\DeclareMathOperator{\im}{im}
\DeclareMathOperator{\var}{var}
\DeclareMathOperator{\Orb}{Orb}
\DeclareMathOperator{\Fix}{Fix}
\DeclareMathOperator{\Stab}{Stab}
\DeclareMathOperator{\cov}{cov}
\DeclareMathOperator*{\esssup}{ess\,sup}
\DeclareMathOperator{\corr}{corr}
\DeclareMathOperator{\lik}{lik}
\DeclareMathOperator*{\argmin}{argmin}
\DeclareMathOperator*{\argmax}{argmax}

\newcommand{\niceline}[2]{%
		\nointerlineskip \vspace{.5\baselineskip}\hspace{\fill}
		{\color{#1}
				\resizebox{0.5\linewidth}{2ex}
				{{%
								{\begin{tikzpicture}
										\node  (C) at (0,0) {};
										\node (D) at (9,0) {};
										\path (C) to [ornament=#2] (D);
										\end{tikzpicture}}}}}%
		\hspace{\fill}
		\par\nointerlineskip \vspace{.5\baselineskip}
}

\definecolor{darkViolet}{HTML}{9400D3}
\newcommand{\sweetline}{%
		\noindent
		\begin{center}
				{\color{darkViolet}
						\resizebox{0.5\linewidth}{1ex}
						{{%
										{\begin{tikzpicture}
												\node  (C) at (0,0) {};
												\node (D) at (9,0) {};
												\path (C) to [ornament=85] (D);
												\end{tikzpicture}}}}}%
		\end{center}
}

\definecolor{remarkPurple}{HTML}{8346FF}
\definecolor{defBlue}{HTML}{0673FF}
\definecolor{exPurple}{HTML}{FF8710}

%THEOREM
\newtcbtheorem[auto counter,number within=section]{theorem}{Theorem}%
{enhanced,colback=white, breakable,frame empty,interior empty,colframe=cyan!50!white, top=8mm,
				coltitle=black,fonttitle=\bfseries,colbacktitle=cyan!15!white,
				borderline={0.5mm}{0mm}{cyan!15!white},
				borderline={0.5mm}{0mm}{cyan!50!white,dashed},
				attach boxed title to top left={yshift=-4mm},
				boxed title style={sharp corners=east,boxrule=1pt},varwidth boxed title}{thm}

%PROPOSITION
\newtcbtheorem[use counter from=theorem]{proposition}{Proposition}%
{enhanced,colback=white,breakable,frame empty,interior empty,colframe=defBlue!75!white, top=8mm,
				coltitle=black,fonttitle=\bfseries,colbacktitle=defBlue!20!white,
				borderline={0.5mm}{0mm}{defBlue!20!white},
				borderline={0.5mm}{0mm}{defBlue!50!white,dashed},
				attach boxed title to top left={yshift=-4mm},
				boxed title style={sharp corners=east,boxrule=1pt},varwidth boxed title}{prop}

%DEFINITION
\newtcbtheorem[use counter from=theorem]{definition}{Definition}%
{enhanced,colback=white,breakable,frame empty,interior empty,colframe=defBlue!75!white, top=8mm,
				coltitle=black,fonttitle=\bfseries,colbacktitle=defBlue!20!white,
				borderline={0.5mm}{0mm}{defBlue!20!white},
				borderline={0.5mm}{0mm}{defBlue!50!white,dashed},
				attach boxed title to top left={yshift=-4mm},
				boxed title style={sharp corners=east,boxrule=1pt},varwidth boxed title}{def}

%COROLLARY
\newtcbtheorem[use counter from=theorem]{corollary}{Corollary}%
{enhanced,colback=white,breakable,frame empty,interior empty,colframe=defBlue!75!white, top=8mm,
				coltitle=black,fonttitle=\bfseries,colbacktitle=defBlue!20!white,
				borderline={0.5mm}{0mm}{defBlue!20!white},
				borderline={0.5mm}{0mm}{defBlue!50!white,dashed},
				attach boxed title to top left={yshift=-4mm},
				boxed title style={sharp corners=east,boxrule=1pt},varwidth boxed title}{cor}

%REMARK
\newtcbtheorem[no counter]{remark}{Remark}%
{detach title, colback=white,enhanced ,breakable,frame empty, interior empty, fonttitle=\bfseries, coltitle=Violet, before upper={\tcbtitle.\quad},
				borderline west={0.5mm}{0mm}{remarkPurple!40!white},
				borderline west={0.5mm}{0mm}{remarkPurple!60!white,dashed}}{remark}

%LEMMA
\makeatletter
\newtcbtheorem[number within = tcb@cnt@theorem]{lemma}{Lemma}%
{enhanced,breakable,colback=white,frame empty,interior empty,colframe=orange!75!white, top=8mm,
				coltitle=black,fonttitle=\bfseries,colbacktitle=orange!20!white,
				borderline={0.5mm}{0mm}{orange!20!white},
				borderline={0.5mm}{0mm}{orange!50!white,dashed},
				attach boxed title to top left={yshift=-4mm},
				boxed title style={sharp corners=east,boxrule=1pt},varwidth boxed title}{lemma}
\makeatother


%PROOF
%%{enhanced,breakable,frame empty,interior empty,colframe=remarkPurple!75!white, top=8mm,
%	coltitle=black,fonttitle=\bfseries,colbacktitle=remarkPurple!20!white,
%	borderline={0.5mm}{0mm}{remarkPurple!20!white},
%	borderline={0.5mm}{0mm}{remarkPurple!50!white,dashed},
%	attach boxed title to top left={yshift=-4mm},
%	boxed title style={sharp corners=east,boxrule=1pt},varwidth boxed title}{prf}


\tcolorboxenvironment{proof}{% amsthm' 
				blanker,breakable,left=5mm,
				before skip=10pt,after skip=10pt,
				borderline west={0.5mm}{0pt}{cyan!40},
				borderline west={0.5mm}{0pt}{remarkPurple!10, dashed}}

%PROBLEM
\newtcbtheorem[auto counter]{problem}{Problem}%
{enhanced,breakable,colback=white,frame empty,interior empty,colframe=cyan!50!white, top=8mm,
				coltitle=black,fonttitle=\bfseries,colbacktitle=cyan!20!white,
				borderline={0.5mm}{0mm}{cyan!20!white},
				borderline={0.5mm}{0mm}{cyan!50!white,dashed},
				attach boxed title to top left={yshift=-4mm},
				boxed title style={sharp corners=east,boxrule=1pt},varwidth boxed title}{prob}

%EXAMPLE
%\newtcbtheorem[use counter from=problem]{example}{Example}%
%{enhanced,breakable,colback=white,frame empty,interior empty,colframe=remarkPurple!50!white, top=8mm,
%		coltitle=black,fonttitle=\bfseries,colbacktitle=remarkPurple!30!white,
%		borderline={0.5mm}{0mm}{remarkPurple!30!white},
%		borderline={0.5mm}{0mm}{remarkPurple!30!white,dashed},
%		attach boxed title to top left={yshift=-4mm},
%		boxed title style={sharp corners=east,boxrule=1pt},varwidth boxed title}{ex}


\newtcbtheorem[use counter from=theorem]{example}{Example}%
{detach title, colback=white,enhanced ,breakable,frame empty, interior empty, fonttitle=\bfseries, coltitle=black, before upper={\tcbtitle.\quad},
		borderline west={0.5mm}{0mm}{remarkPurple!30!white},
		borderline ={0.5mm}{0mm}{remarkPurple!30!white}}{example}

%SOLUTION
\newtcbtheorem[no counter]{solution}{Solution}%
{enhanced,breakable,colback=white,frame empty,interior empty,colframe=green!75!white, top=8mm,
				coltitle=black,fonttitle=\bfseries,colbacktitle=green!20!white,
				borderline={0.5mm}{0mm}{green!20!white},
				borderline={0.5mm}{0mm}{green!50!white,dashed},
				attach boxed title to top left={yshift=-4mm},
				boxed title style={sharp corners=east,boxrule=1pt},varwidth boxed title}{sol}
\definecolor{realPurple}{HTML}{AA05F9}
\definecolor{gray}{rgb}{0.5,0.5,0.5}
\definecolor{dkgreen}{rgb}{0,0.6,0}
\definecolor{mauve}{rgb}{0.58,0,0.82}

\lstset{frame=tb,
				style=Matlab-editor,
				language=C,
				aboveskip=3mm,
				belowskip=3mm,
				xleftmargin=3mm,
				showstringspaces=false,
				columns=flexible,
				frame=none,
				basicstyle={\small\ttfamily},
				numberstyle=\tiny\color{gray},
				keywordstyle=\color{blue},
				commentstyle=\color{dkgreen},
				stringstyle=\color{mauve},
				breaklines=true,
				breakatwhitespace=true,
				mlshowsectionrules = true,
				tabsize=3,
				backgroundcolor=\color{cyan!5}
}

\newcommand\mmybox[2][fill=cyan!20]{%
    \tikz[baseline]\node[%
        inner ysep=0pt, 
        inner xsep=2pt, 
        anchor=text, 
        rectangle, 
        rounded corners=1mm,
        #1] {\strut#2};%
}


\def\changemargin#1#2{\list{}{\rightmargin#2\leftmargin#1}\item[]}
\let\endchangemargin=\endlist

\linespread{1.4}



% MAIN DOC
\begin{document}

\title{Dynamic Programming}
\author{Adam Yang}
% \date{\today}
\maketitle


\begin{definition*}{High Level Understanding}
    One way to think about it: avoid duplicating work in backtracking solutions (when it is possible).

    Closely related to recursive thinking: What smaller problems do I need to have solved to be able to now solve a bigger problem?

    Contrast to Greedy: What is our first step? Commit to it, then solve a new problem.
\end{definition*}

\subsection*{Warm-Up}

\begin{problem*}{}
    Find the n-th Fib Number
\end{problem*}
\begin{solution*}{recursive, not optimal}
            Let $f(n)$ be the n-th fib number, we have
        
	\[f(0)=f(1)=1\]
        \[n\geq 2: f(n)=f(n-1)+f(n-2)\]
        \[T(n)=T(n-1)+T(n-2)+O(1)\]

        So $T(n)\geq f(n)$, which grows roughly as Golden-Ratio.
        
        Golden Ratio is the $x$ solving $x^2=x+1, x\approx 1.618$. The running time is exponential.
\end{solution*}
		% \renewcommand{\qedsymbol}{} % hide the QED square



 \begin{proof}[Improvement]{}
        		\renewcommand{\qedsymbol}{} % hide the QED square
                This recalculates $f(n-2)$ twice, and as a result, recalculates $f(n-3)$ for more often, and so on. We should instead store the result once we have computed it. At that point, we get rid of recursion:

        \qquad \texttt{int Fib[n+1];}\texttt{}
        
        \qquad \texttt{Fib[0] = Fib[1] = 1;}
        
        \qquad \texttt{for(int i=2; i $\leq$ n; ++i) Fib[i]=Fib[i-1]+Fib[i-2];}

        This makes the running time to $O(n)$
        \end{proof}
% unnumbered theorem/proposition/definition/example/etc: add *

\section*{Problem}
\begin{problem*}{Weighted Interval Selection}
    Same as our original selection problem:
    $n$ intervals with start times s(i), finish time f(i), and \textbf{*weights*} w(i).

    \textbf{Goal:} select a set S of intervals where no two intervals overlap, and so that the total weight of selected intervals is maximized.

    \textbf{Application:} weights could be importance/interest of events between which we need to choose
\end{problem*}

\begin{proof}[Analysis]{}
		\renewcommand{\qedsymbol}{} % hide the QED square
        Greedy gives us suboptimal solution. \textbf{No greedy rule will work}
\end{proof}
\begin{proof}[Key Insight:]{}
		\renewcommand{\qedsymbol}{} % hide the QED square
        Exhaustive search works, but the running time is $2^n$. 

        If we have intervals 1, ..., n available, the optimum solution for those intervals does one of the following (with respect to interval $i$):

        \qquad 1) Not include $i$, and get the optimum solution for all intervals except $i$.
        
        \qquad 2) Include $i$, and get the optimum solution for all intervals except those intersecting $i$, plus the weight that it gets from $i$ itself.

        \qquad Among these, the optimum solution (being optimal) chooses the better one.

        What makes this tick is that once we commit to including $i$, or to not including $i$, we compute/use the \textbf{*optimum*} solution for the resulting subproblem. \textbf{This is the key property for DP:} Any optimal soution must contain optimal solutions to the subproblems contained in it. We have to define the subproblems to that this is true, in order to be able to apply DP.
\end{proof}

\begin{theorem*}{}
    For DP, we often analyze relatively easy (or even trivial) properties of the optimum solution, and exploit them to derive an efficient algorithm
\end{theorem*}
\begin{solution*}{}
    Define OPT(S):= maximum weight that can be obtained from intervals in S.

    \begin{center}
            \textbf{[Find the max from including interval i or not including i]}
    \end{center}
    
    OPT(S) = max( OPT(s\textbackslash \{i\}), w(i)+OPT(S \textbackslash \{everything intersecting i\}) )

    \textbf{Base Case:} OPT(\{\}) = 0

    \begin{proof}{This is not optimal if we do it recursively}
    
        Implementing this \textbf{recursively} is just exhaustive search with backtracking, and runs in about $2^n$ times.

        If we store all OPT(S) values in a table to avoid recomputing, with the current approach, we still take $2^n$ and the able size is $2^n$
    \end{proof}

    Our power is that we can choose the interval number i that we branch over. By a careful choice, we can make sure that only polynomially many (in fact: linearly many) subproblems need to solved.

    Assume from now on that intervals are sorted by non-decreasing finish time f(i). Subproblems of the form \{1,2,...,k\} for k = 0,...,n.

    We define the following notations
    
    OPT(k):= optimum weight for the subproblems \{1,2,...,k\}
    
    t(k):= the largest index j such that $f(j) < s(k)$ [the latest interval that does not intersect interval k.]
    
    Recurrence Relation:
    \[OPT(k) = max(OPT(k-1), w(k)+OPT(t(k)))\]

    \textbf{Base Case:} OPT(0) = 0
    
    \textbf{Justification:} The opt solution for intervals 1,..,k either doesn't contain k (in which case it is the OPT(k-1)) or it does contain k, and includes the opt solution for the still available intervals

    Final answer we're interested in: OPT(n)

    \textbf{Implementation:}

    \qquad \texttt{1. Sort intervals by non-decreasing f(i)}
    
    \qquad \texttt{2. Precompute all of the t[k];}
    
    \qquad \texttt{3. a[0] = 0;}
    
    \qquad \texttt{4. for(int i = 1; i <= n; ++i) a[i] = max(a[i-1], w[i]+a[t[i]]);}
    
    \qquad \texttt{5. return a[n];}
\end{solution*}

To output the solution itself:
k=n; while(k>0) if (k yellow) {output k; k = t[k];} else k--;

% \begin{proof}[Problem Insight]{}
% \renewcommand{\qedsymbol}{}
%         what do we know for sure about the opt solution set S, and on particular interval?

        
% \end{proof}

% END MAIN Doc
\end{document}
